\documentclass{article}
\usepackage[T1]{fontenc}
\usepackage{lmodern}
\usepackage{graphicx}
\usepackage{amsmath,amssymb}
\usepackage[a4paper,margin=2cm]{geometry}
\usepackage[absolute,overlay]{textpos}
\setlength{\TPHorizModule}{1mm}
\setlength{\TPVertModule}{1mm}
\usepackage[indonesian]{babel}
\usepackage{hyperref}
\setlength{\emergencystretch}{2em}
% unit: Halaman 40

\begin{document}

% === Halaman 40 (python/native) ===

40

7.  Cipher

 -  Algoritma untuk enkripsi dan dekripsi pesan 

 -  Aturan (rule) untuk enchipering dan  dechipering, atau berupa

    fungsi matematika yang digunakan untuk enkripsi dan dekripsi  pesan.

   Contoh:   Enkripsi: Geser tiga huruf ke kanan     (rule)


         Dekripsi: Geser tiga huruf ke kiri          (rule) 


         E(p) = (p + k) mod 26                 (fungsi matematika)


         D(c) = (c – k) mod 26 

          (fungsi matematika)


   - Classical cipher: Caesar cipher, Vigenere cipher, Playfair Cipher, Enigma cipher, dll

   - Modern cipher: DES, AES, Blowfish, Serpent, RSA, ElGamal, RC4, RC5, A5, dll

\end{document}
